\bookchapter{Chapitre 32 : Manuel d'urgence de la sécurité mobile}

\chapterintro{Ce chapitre est un « plan d'urgence ». Nous espérons que vous n'en aurez jamais besoin, mais en cas de perte, de vol ou d'arnaque, il vous aidera à limiter les dégâts ; lors de séismes ou de fortes pluies, le téléphone sert aussi à recevoir les alertes, à se localiser et à rester en contact. Lisez-le maintenant, puis capturez ou imprimez les étapes clés et gardez-les chez vous — au moment critique, votre téléphone n'est peut-être plus entre vos mains, ou sans batterie ni réseau.}

\booksection{Scénario zéro : séisme / intempéries — comment le téléphone vous aide à survivre}

Ici, il s'agit de protéger les personnes, pas les comptes. Pour les réglages d'alerte, voir le chapitre 13 « Alertes météo, catastrophes naturelles et séismes » ; nous ne décrivons ici que la conduite à tenir une fois l'alarme retentie.

\textbf{Alarme sismique : se protéger d'abord, évacuer ensuite — ne filmez pas}\par

\begin{enumerate}
  \item Protégez-vous sur place : s'abaisser, se couvrir, s'agripper (sous une table, dans un angle porteur, près d'un meuble solide) ; protégez tête et cou ; loin des vitres, murs extérieurs et lustres.
  \item Ne fuyez pas dans l'escalier : la secousse principale dure souvent une dizaine à quelques dizaines de secondes ; la cage d'escalier peut être plus dangereuse. L'ascenseur est formellement interdit.
  \item Une fois la secousse passée, évacuez en ordre : par les escaliers, attention aux répliques et aux chutes d'objets ; gagnez un espace dégagé, loin des poteaux, panneaux et bâtiments fragiles.
  \item Utilisez le téléphone au bon moment : un court SMS à vos contacts d'urgence (« Je suis à XX, en sécurité/blessé, j'ai besoin d'aide ») consomme moins qu'un long appel ; privilégiez le texte si possible.
  \item Sans réseau : voir le chapitre 7 SOS satellite / messages satellite (si le modèle le permet) ; ou montez en hauteur, dans un endroit ouvert ; dans certains cas, la talkie-walkie Bluetooth hors ligne aide aussi.
  \item Coupez le superflu : mode avion activé/désactivé par cycles, Bluetooth et scan désactivés, luminosité basse — plusieurs heures de batterie en plus. Gardez la dernière barre pour rassurer vos proches et appeler les secours.
\end{enumerate}

\begin{tcbwarn}{Pluies torrentielles / typhon / crues : agissez dès l'alerte}
\begin{itemize}
  \item Alerte rouge : limitez les sorties, éloignez-vous des sous-sols, ponts bas, lits de rivière et pentes raides.
  \item Bloqué dehors : gagnez les étages supérieurs d'un bâtiment solide, ne traversez pas une zone inondée (bouches d'égout absentes, risque d'électrocution) ; en appelant le 119/110, précisez les repères et le niveau de l'eau.
  \item Avant de partir : vérifiez les risques via l'application Météo et les comptes d'alerte locaux ; en long trajet, évitez les routes de montagne le long de l'eau ou des falaises.
  \item Famille : convenez que « en cas de perte de contact, on se fie à la dernière localisation WeChat/SMS » pour éviter d'épuiser la batterie en appels répétés.
\end{itemize}

\end{tcbwarn}

\begin{tcbpractice}{4 réflexes à prendre dès maintenant}
\begin{enumerate}
  \item Activez l'alerte sismique système ou une application d'alerte fiable, et autorisez-la à percer le mode silencieux.
  \item Épinglez 2 à 3 « contacts d'urgence » en tête de liste ; affichez-les sur l'écran de verrouillage (fiche médicale iOS/Android ou informations d'écran).
  \item Gardez une batterie externe à la maison et dans le sac ; chargez-la avant un long déplacement et emportez le câble.
  \item Mémorisez au moins un numéro appelable sans les contacts (le mobile d'un proche) et notez-le sur un papier dans le portefeuille.
\end{enumerate}

\end{tcbpractice}

\booksection{Scénario un : téléphone perdu/volé — les 30 minutes d'or}

Les 30 premières minutes après la perte sont « d'or » : bien agir protège votre argent et votre vie privée. Suivez l'ordre ci-dessous sans sauter d'étape :

\textbf{Procédure standard des 30 minutes}\par

\begin{enumerate}
  \item Appelez votre propre numéro depuis un autre téléphone (minute 1) : il a peut-être juste été laissé à proximité ; une personne honnête pourra vous le rendre.
  \item Verrouillez à distance (minutes 2-5) : depuis un autre appareil, connectez-vous au service de localisation : iPhone : iCloud.com > Localiser > votre appareil > Marquer comme perdu
  \item Android selon la marque : Huawei cloud.huawei.com / Xiaomi i.mi.com / OPPO cloud.oppo.com, etc.
  \item Faites opposition sur la carte SIM (minutes 5-10) : appelez le service client de l'opérateur (China Mobile 10086 / China Unicom 10010 / China Telecom 10000) et demandez la suspension de la SIM. Étape la plus critique — avec votre SIM, l'escroc reçoit les codes et réinitialise tous vos mots de passe.
  \item Gélez les comptes de paiement (minutes 10-20) : Alipay : opposition au 95188, ou depuis un autre téléphone Alipay > Moi > Réglages > Centre de sécurité > Gel d'urgence
  \item WeChat : après dépôt de plainte (110), contactez le support WeChat, ou Centre de sécurité WeChat > Geler le compte
  \item Application bancaire : appelez le numéro au dos de la carte et demandez le gel de la banque en ligne
  \item Prévenez vos proches (minutes 20-25) : via Moments ou par un membre de la famille, informez que votre téléphone est perdu, pour empêcher l'escroc de se faire passer pour vous.
  \item Portez plainte (minutes 25-30) : appelez le 110 ou rendez-vous au commissariat le plus proche et obtenez un récépissé. Il sera nécessaire pour l'assurance et les litiges bancaires.
\end{enumerate}

\booksection{Scénario deux : vous avez été escroqué — 72 heures d'or}

En cas d'arnaque, le plus important est de garder la tête froide. Des décisions prises dans la panique entraînent une seconde perte. Procédure :

\begin{enumerate}
  \item Cessez immédiatement tout virement : quoi que dise l'escroc, ne transférez plus un centime. « Encore un versement de garantie et on vous rembourse » est l'arnaque classique de la deuxième vague.
  \item Conservez toutes les preuves : ne supprimez ni l'historique de discussion ni les journaux d'appels ! Capturez les conversations, les virements et les coordonnées du destinataire. Ce sont les preuves clés pour la plainte et le recouvrement.
  \item Composez immédiatement le 96110 : ligne nationale anti-fraude. L'opérateur vous guidera pour un arrêt de paiement d'urgence — si le virement est récent (souvent moins de 24 h), la banque peut geler le compte destinataire et intercepter les fonds.
  \item Portez plainte au 110 : avec téléphone, pièce d'identité et justificatifs de virement, rendez-vous au commissariat pour déposer plainte.
  \item Contactez la banque/la plateforme de paiement : signalez l'arnaque et demandez le gel des comptes concernés. En cas de fraude sur carte de crédit, demandez la procédure de refus de paiement.
  \item Changez tous les mots de passe : depuis un autre appareil sûr, modifiez les mots de passe Apple ID/compte de marque, WeChat, Alipay et applications bancaires, et activez la double authentification.
\end{enumerate}

\booksection{Scénario trois : virus/cheval de Troie sur le téléphone — que faire ?}

Si des publicités apparaissent sans raison, si des applications bizarres s'installent seules, si la facture explose ou si les données partent anormalement, un logiciel malveillant est probablement présent :

\begin{itemize}
  \item Étape 1 : mode sans échec — sur Android, maintenez le bouton d'alimentation pour éteindre, rallumez, puis maintenez la touche volume bas dès la vibration jusqu'à l'apparition de « Mode sans échec » en bas à l'écran. Les applications tierces sont alors désactivées, ce qui facilite le diagnostic.
  \item Étape 2 : désinstallez les applications suspectes — en mode sans échec, Réglages > Gestion des applications, triez par date et supprimez les applications inconnues récemment installées.
  \item Étape 3 : scan de sécurité — installez « Tencent Mobile Manager » ou « 360 Mobile Guard » et lancez une analyse complète.
  \item Étape 4 : réinitialisation d'usine (solution ultime) — si rien ne fonctionne, sauvegardez les données importantes puis réinitialisez. Attention : tout est effacé, sauvegardez d'abord photos, contacts et documents.
  \item iPhone : les infections sont rares — le bac à sable d'iOS empêche la plupart des malwares. En cas d'anomalie, il s'agit souvent d'un bug d'application ou d'un réglage. Une réinitialisation résout 99 \% des problèmes.
\end{itemize}

\booksection{Scénario quatre : compte piraté — reprendre vite le contrôle}

\begin{enumerate}
  \item Changez immédiatement le mot de passe : si vous pouvez encore vous connecter, faites-le tout de suite. Le nouveau doit être totalement différent de l'ancien.
  \item Déconnectez les autres appareils : dans les réglages de sécurité, « Gestion des appareils connectés » et déconnectez tous les appareils inconnus.
  \item Vérifiez l'historique de connexion : sur 7 jours, une connexion depuis un autre lieu signe un piratage.
  \item Activez la double authentification : si ce n'est pas encore fait, faites-le maintenant. C'est la protection la plus efficace contre une récidive.
  \item Vérifiez les informations liées : numéro et e-mail ont-ils été modifiés ? Si oui, utilisez la procédure de récupération du compte.
  \item Vérifiez les paiements : pour WeChat/Alipay, contrôlez les paiements sans mot de passe et les prélèvements automatiques, ainsi que les cartes bancaires ajoutées à des comptes inconnus.
\end{enumerate}

\booksection{Scénario cinq : message de chantage — n'ayez pas peur, ne payez pas}

Deux escroqueries courantes : (1) « J'ai une vidéo/photo compromettante de vous ; payez ou j'envoie à tous vos contacts » ; (2) « Votre téléphone est piraté ; payez ou il sera verrouillé ». Face à cela :

\begin{itemize}
  \item 99,9 \% d'arnaques : ils n'ont aucune donnée privée vous concernant. En envoyant le même message à des dizaines de milliers de personnes, certaines paient par peur.
  \item Ne répondez pas, ne payez pas : répondre confirme que « ce numéro est actif » et vous classera comme cible de chantage ; les harcèlements se multiplieront.
  \item Capturez l'écran en preuve, puis bloquez l'expéditeur.
  \item Changez les mots de passe : par précaution, mettez à jour ceux de vos comptes principaux.
  \item Si une vraie fuite de vie privée est possible : portez plainte. La police a l'habitude de ce type d'affaires.
\end{itemize}

\booksection{Scénario six : arnaques par deepfake vocal/visage IA — l'image n'est plus une preuve}

Depuis 2024, le changement de visage et de voix par IA permet aux escrocs de « devenir » un proche, un supérieur ou un ami en visio. « Maman, j'ai un problème, envoie vite de l'argent » — la personne en vidéo ressemble et parle exactement comme votre fils, mais c'est un faux. C'est l'une des arnaques les plus difficiles à détecter.

\textbf{Étape 1 : repérer les failles — quatre indices courants d'un deepfake}\par

\begin{itemize}
  \item Clignements anormaux : une personne cligne 15 à 20 fois par minute ; un visage IA cligne trop peu ou de façon mécanique.
  \item Lèvres désynchronisées : un décalage de quelques dixièmes de seconde entre le mouvement des lèvres et le son.
  \item Contours flous : bords du visage, oreilles, cheveux souvent granuleux ou scintillants.
  \item Lumière incohérente : reflets et ombres du visage ne correspondent pas à l'éclairage du fond ; de profil, les traits « tournent » de façon bizarre.
\end{itemize}

\begin{tcbinfo}{Étape 2 : trois vérifications}
\begin{itemize}
  \item Demandez un geste précis : tourner la tête, se couvrir le visage, coucher la bouche avec la main puis l'ouvrir. Un deepfake « craque » plus facilement quand le visage est masqué ou en forte rotation.
  \item Posez une question que lui seul connaît : inventez-en une, par exemple « dans quel hôpital a-t-on emmené le chien la dernière fois ? » — un escroc bafouillera.
  \item Rappelez un numéro connu : raccrochez la visio et composez le numéro enregistré dans vos contacts. L'escroc ne contrôle pas le vrai numéro — l'arnaque s'effondre immédiatement.
\end{itemize}

\end{tcbinfo}

\begin{tcbwarn}{Étape 3 : procédure de plainte après confirmation}
\begin{enumerate}
  \item Stoppez tout transfert : même sous l'urgence, faites une pause ; quiconque insiste pour « virement rapide, ne raccrochez pas » est presque toujours un escroc.
  \item Conservez toutes les preuves : enregistrement d'écran, captures de l'appel vidéo, conversations, comptes et justificatifs — matériaux essentiels pour l'expertise technique.
  \item Appelez le 96110 : indiquez qu'il s'agit d'une arnaque « deepfake visage/voix IA » ; l'opérateur lancera l'arrêt de paiement d'urgence.
  \item Portez plainte au 110 : apportez toutes les preuves et signalez l'usage de l'IA, pour faciliter la qualification judiciaire.
  \item Prévenez vos proches : les escrocs peuvent se servir de votre vidéo falsifiée pour tromper d'autres personnes ; postez un avertissement dans Moments ou le groupe familial.
\end{enumerate}

\end{tcbwarn}

\booksection{Scénario sept : attaque SIM swap — votre numéro « volé »}

Peu connue, mais dévastatrice. Le SIM swap consiste pour un attaquant à se faire passer pour vous en boutique et à obtenir une nouvelle SIM avec vos identifiants. Votre téléphone passe alors « hors service » tandis que la nouvelle carte dans son appareil reçoit tous vos codes — Alipay, banque, WeChat : toutes les portes sont ouvertes.

\begin{tcbinfo}{Comment le détecter ? Trois signaux d'alerte}
\begin{itemize}
  \item Le téléphone affiche soudain « Pas de service » ou « Pas de carte SIM » sans que vous y ayez touché, même après redémarrage — c'est souvent le moment du détournement.
  \item Vous recevez un SMS de l'opérateur confirmant un « remplacement de carte » ou une « opération réussie » alors que vous n'avez rien demandé.
  \item Vous ne recevez plus les codes, ni même les autres SMS — les messages sont détournés vers la carte de l'attaquant.
\end{itemize}

\end{tcbinfo}

\begin{tcbwarn}{Protection — verrouillez aussi la carte SIM}
\begin{itemize}
  \item Définissez un code PIN SIM : chemin Réglages > Téléphone > Code PIN de la carte SIM. Après activation, chaque redémarrage ou changement de carte exige le PIN. Même avec la carte, l'attaquant ne peut rien en faire.
  \item Activez la double vérification de l'opérateur : dans les applications China Mobile/Unicom/Telecom, « Centre de sécurité » ou « Vérification des opérations », activez l'alerte SMS en cas de remplacement de carte par un tiers et le second mot de passe. En boutique, une vérification d'identité supplémentaire sera alors exigée.
  \item Pour les comptes sensibles, évitez le seul SMS : banque, Apple ID, etc. — privilégiez une application d'authentification (Google Authenticator, Microsoft Authenticator) ou une clé matérielle.
  \item Filigranez les copies de pièce d'identité : indiquez « Réservé à l'usage de XX » — les attaquants s'appuient souvent sur des identifiants fuités pour usurper votre identité.
\end{itemize}

\end{tcbwarn}

\textbf{Après un détournement — la vitesse est tout}\par

\begin{enumerate}
  \item Appelez immédiatement le service client pour opposition : China Mobile 10086 / China Unicom 10010 / China Telecom 10000 ; suspendez la ligne et bloquez le remplacement de carte.
  \item Connectez-vous depuis un autre appareil : changez sans tarder les mots de passe banque, Alipay, WeChat, Apple ID et déconnectez les appareils inconnus.
  \item Gélez les comptes de paiement : Alipay 95188, Centre de sécurité WeChat, service client bancaire pour la banque en ligne.
  \item En boutique avec votre pièce d'identité : obtenez une nouvelle carte dans un magasin officiel et recouvrez l'usage de votre numéro.
  \item Portez plainte : le SIM swap constitue un « délit d'atteinte aux informations personnelles des citoyens » ; apportez les preuves (délai de coupure, préjudice financier, enregistrement du remplacement).
  \item Audit complet : une fois le numéro récupéré, vérifiez connexions et mouvements de fonds récents pour toute trace d'utilisation frauduleuse.
\end{enumerate}

\begin{tcbwarn}{Mieux vaut prévenir — 5 réflexes quotidiens}
\begin{enumerate}
  \item Code de verrouillage + double authentification, les deux activés
  \item « Localiser mon téléphone » bien activé
  \item Code PIN SIM défini (Réglages > Téléphone > Code PIN de la carte SIM) — si le téléphone est perdu, personne n'utilise votre carte
  \item N'installez pas d'applications d'origine inconnue, ne cliquez pas sur des liens suspects
  \item Sauvegardez régulièrement les données importantes (photos, contacts, discussions)
\end{enumerate}

\end{tcbwarn}

\begin{tcbpractice}{À faire : créez votre « carte de contacts d'urgence »}
\begin{enumerate}
  \item Sur un papier, notez : service client des opérateurs (China Mobile 10086 / China Unicom 10010 / China Telecom 10000), service client bancaire, ligne anti-fraude 96110, police 110, téléphone de la famille.
  \item Notez l'e-mail de votre Apple ID/compte de marque et un contact de secours.
  \item Gardez ce papier dans le portefeuille ou sur le réfrigérateur — pas dans le téléphone, qui peut ne pas être là au moment critique.
  \item Vérifiez dès maintenant : « Localiser mon téléphone » est-il actif ? Le code PIN SIM est-il défini ?
\end{enumerate}

\end{tcbpractice}

\begin{tcbsummary}{Résumé du chapitre}
\begin{itemize}
  \item Téléphone perdu, 30 minutes d'or : appeler depuis un autre téléphone verrouiller à distance opposition SIM geler les paiements prévenir les proches porter plainte
  \item Après une arnaque, 72 heures : stopper les virements conserver les preuves appeler le 96110 porter plainte contacter la banque changer les mots de passe
  \item Téléphone infecté : mode sans échec désinstaller les applications suspectes scan de sécurité réinitialisation (ultime)
  \item Compte piraté : changer le mot de passe déconnecter les appareils vérifier l'historique activer la 2FA vérifier les liens vérifier les paiements
  \item Messages de chantage : 99,9 \% d'arnaques — ne pas répondre, ne pas payer, capturer et bloquer
  \item Deepfake visage/voix : clignements, lèvres, contours, lumière ; demander un geste, poser une question privée, rappeler un numéro connu
  \item SIM swap : « Pas de service » soudain = alerte ; code PIN SIM, double vérification opérateur, authentificateur pour les comptes importants
  \item 5 préventions quotidiennes : verrouillage + 2FA, localisation du téléphone, PIN SIM, pas d'applications douteuses, sauvegardes régulières
\end{itemize}

\end{tcbsummary}
